\Needspace{10\baselineskip}
\centerline{\textsc{Gauß an Bode.}}

\medskip
Astronomisches Jahrbuch für 1808. S. 187---190. Berlin 1808.

\medskip
\begin{flushright}
Braunschweig, 30. September 1804.
\end{flushright}

Ew. --- erlaube ich mir, hier alle meine bisherigen Beobachtungen der Juno vorzulegen. Ich habe sie zwar nur mit dem Kreismikrometer, aber mit größtmöglicher Sorgfalt angestellt. Sie gehen sämtlich aus sehr zahlreichen Vergleichen hervor, teils mit Sternen aus \textsc{Piazzi}s Katalog, teils mit Sternen aus der Hist. Céleste, die ich nach diesem Katalog sorgfältig reduziert habe. Bis auf die ersten 3, die mit einem schlechteren Instrument gemacht wurden, habe ich alle mit einem sehr guten Spiegelteleskop von \textsc{Short} angestellt.

\begin{center}
\begingroup\fontsize{10}{12}\selectfont\setlength{\tabcolsep}{3pt}\renewcommand{\arraystretch}{1.12}
\begin{tabular}{r|rrr|rrr|rrr}
1804 & \multicolumn{3}{c|}{Mittlere Zeit} & \multicolumn{6}{c}{Scheinbare} \\ \cline{5-10}
 & \multicolumn{3}{c|}{} & \multicolumn{3}{c|}{Rektaszension} & \multicolumn{3}{c}{Südliche Deklination} \\ \hline
12. Sept. & $10^u$ & $35^m$ & $2^s$ & $0^\circ$ & $54'$ & $31''$ & $1^\circ$ & $38'$ & $15''$ \\
13 & $9$ & $41$ & $32$ &  & $45$ & $24$ &  & $50$ & $59$ \\
14 & $11$ & $38$ & $49$ &  & $35$ & $37$ & $2$ & $5$ & $1$ \\
15 & $10$ & $16$ & $17$ &  & $26$ & $53$ &  & $17$ & $35$ \\
16 & $10$ & $37$ & $4$ &  & $17$ & $17$ &  & $31$ & $20$ \\
17 & $11$ & $28$ & $59$ &  & $7$ & $23$ &  & $44$ & $29$ \\
\end{tabular}\endgroup
\end{center}
\begin{center}
\begingroup\fontsize{10}{12}\selectfont\setlength{\tabcolsep}{3pt}\renewcommand{\arraystretch}{1.12}
\begin{tabular}{r|rrr|rrr|rrr}
1804 & \multicolumn{3}{c|}{Mittlere Zeit} & \multicolumn{6}{c}{Scheinbare} \\ \cline{5-10}
 & \multicolumn{3}{c|}{} & \multicolumn{3}{c|}{Rektaszension} & \multicolumn{3}{c}{Südliche Deklination} \\ \hline
18. Sept. & $11^u$ & $22^m$ & $16^s$ & $359^\circ$ & $57'$ & $26''$ &  & $58'$ & $3''$ \\
21 & $10$ & $24$ & $52$ &  & $28$ & $6$ & $3$ & $38$ & $0$ \\
24 & $10$ & $1$ & $2$ & $358$ & $57$ & $53$ & $4$ & $18$ & $6$ \\
25 & $8$ & $44$ & $25$ &  & $48$ & $12$ &  & $30$ & $44$ \\
27 & $10$ & $20$ & $29$ &  & $27$ & $20$ &  & $57$ & $47$ \\
28 & $8$ & $29$ & $4$ &  & $18$ & $20$ & $5$ & $10$ & $22$ \\
\end{tabular}\endgroup
\end{center}
\SourceLeaf{254}%

Teils anhand dieser Beobachtungen, teils anhand der Beobachtungen unseres gemeinsamen Freundes, Herrn Dr. \textsc{Olbers}, und einiger Beobachtungen des Freiherrn \textsc{von Zach} habe ich bereits eine erste Annäherung an die Elemente versucht, die ich die Ehre habe, Ihnen hier mitzuteilen. Ich halte die Planetennatur der Juno nun bereits \emph{aus den Beobachtungen} für erwiesen und glaube hoffen zu dürfen, dass auch die Elemente schon ausreichen, um eine Vorstellung von den ungefähren Dimensionen der Bahn zu geben,

\begin{center}
I\textsuperscript{te} Elemente der Junobahn.
\end{center}

\begin{center}
\begingroup\fontsize{8}{10}\selectfont\setlength{\tabcolsep}{3pt}
\begin{tabular}{lr}
Epoche 5. September 1804, $0^u$ mittlere Zeit im Seeberger Meridian . & $20^\circ\ 38'\ 56''$ \\
Tägliche mittlere Bewegung . . . . . . . . . . & $779''80$ \\
Aphel . . . . . . . . . . . . . . . . & $239^\circ\ 14'\ 2''$ \\
Log. der Halbachse . . . . . . . . . . . . . & $0.438682$ \\
Exzentrizität . . . . . . . . . . . . . . . & $0.287359$ \\
Aufsteigender Knoten . . . . . . . . . . . . . . & $171^\circ\ 15'\ 35''$ \\
Neigung . . . . . . . . . . . . . . . . . . & $13\ 34\ 59$ \\
\end{tabular}\endgroup
\end{center}

Da die Beobachtungen bisher nur einen so kurzen Zeitraum umfassen, werden diese Elemente natürlich noch beträchtlicher Verbesserungen bedürfen. Insbesondere ist die Unsicherheit der mittleren Bewegung noch viel größer als der kleine Unterschied zur mittleren Bewegung von $\CeresSymbol$ und $\PallasSymbol$. Daher kann es durchaus sein, dass alle 3 der Gleichheit noch viel näher kommen.

Um weiter entfernten Astronomen, denen es vielleicht noch nicht gelungen ist, den Planeten aufzufinden, die Suche nach Kräften zu erleichtern, habe ich aus den obigen Elementen eine kleine Ephemeride berechnet. Sie wird seinen Lauf hoffentlich mehrere Wochen lang genau genug angeben, um ihn aufzufinden.

\Needspace{26\baselineskip}
\begin{flushright}
Braunschweig, 12. November 1804.
\end{flushright}

Indem ich Ew.--- für die Beobachtungen, die Sie mir mitgeteilt haben, meinen verbindlichsten Dank ausspreche, berichte ich Ihnen zugleich weiter über die neuesten Ergebnisse meiner Untersuchungen zur Bahn der Juno, obwohl sie bereits vor 14 Tagen berechnet wurden und sich daher bald weiter verbessern lassen werden. Hier die

\begin{center}
III\textsuperscript{ten} Elemente der Juno.
\end{center}

\begin{center}
\begingroup\fontsize{9}{11}\selectfont\setlength{\tabcolsep}{3pt}
\begin{tabular}{lr}
Epoche 1805 im Seeberger Meridian . . . . . . . . & $43^\circ\ 21'\ 2''$ \\
Aphel . . . . . . . . . . . . . . . . . & $233\ 56\ 6$ \\
Neigung . . . . . . . . . . . . . . . . . . & $12\ 52\ 48$ \\
Logarithmus der Halbachse . . . . . . . . . . . . & $0.426699$ \\
Tägliche Bewegung . . . . . . . . . . . . . . & $812''754$ \\
Knoten . . . . . . . . . . . . . . . . . . & $171^\circ\ 0'\ 0''$ \\
Exzentrizität . . . . . . . . . . . . . . . . & $0.263182$ \\
\end{tabular}\endgroup
\end{center}

Dass die mittlere Bewegung der Juno der von $\CeresSymbol$ und $\PallasSymbol$ nicht gleich, sondern größer ist, sehe ich bereits als erwiesen an.

Bei der Berechnung der III\textsuperscript{ten} Elemente lagen mir die Beobachtungen von Herrn \textsc{von Zach} bis zum 21. Oktober vor, mit denen sie bis auf ein paar Sekunden übereinstimmen.
Die Beobachtungen von Dr. \textsc{Maskelyne} liegen mir bisher bis zum 17\textsuperscript{ten} vor. Nach den späteren Beobachtungen steht jedoch fest, dass diese Elemente die Rektaszension jetzt zu groß angeben. Über die Deklination kann ich aber noch nichts entscheiden. Hier der Vergleich mit den neuesten Beobachtungen.
\SourceLeaf{255}%


\begin{center}\begin{tabular}{r|cc|c|cc}
 & \multicolumn{2}{c|}{Hr. Prof. \textsc{Bode}} & & \multicolumn{2}{c}{Hr. Dr. \textsc{Olbers}} \\
 & Rekt. & Dekl. & & Rekt. & Dekl. \\
\hline
20. Okt. & $+33''$ & $+7''$ & 27. Okt. & $+21''$ & $+15''$ \\
2. Nov. & $+48$ & $+5$ & 30 & $+34$ & $+12$ \\
5 & $+46$ & $-15$ & 2. Nov. & $+33$ & $-1$ \\
 & & & 5 & $+39$ & $+25$ \\
 & & & 6 & $+37$ & $+32$ \\
\end{tabular}\end{center}

Am 20. Dezember kommen $\CeresSymbol$ und $\JunoSymbol$ einander so nahe, dass man sie gleichzeitig im Fernrohr beobachten kann.


